\documentclass[aps,preprint]{revtex4}%
\usepackage{amssymb}
\usepackage{amsfonts}
\usepackage{amsmath}
\usepackage{graphicx}%
\setcounter{MaxMatrixCols}{30}
%TCIDATA{OutputFilter=latex2.dll}
%TCIDATA{Version=5.50.0.2960}
%TCIDATA{CSTFile=revtex4.cst}
%TCIDATA{Created=Wednesday, August 21, 2019 13:15:11}
%TCIDATA{LastRevised=Wednesday, September 25, 2019 20:26:33}
%TCIDATA{<META NAME="GraphicsSave" CONTENT="32">}
%TCIDATA{<META NAME="SaveForMode" CONTENT="1">}
%TCIDATA{BibliographyScheme=Manual}
%TCIDATA{<META NAME="DocumentShell" CONTENT="Articles\SW\REVTeX 4">}
%BeginMSIPreambleData
\providecommand{\U}[1]{\protect\rule{.1in}{.1in}}
%EndMSIPreambleData
\newtheorem{theorem}{Theorem}
\newtheorem{acknowledgement}[theorem]{Acknowledgement}
\newtheorem{algorithm}[theorem]{Algorithm}
\newtheorem{axiom}[theorem]{Axiom}
\newtheorem{claim}[theorem]{Claim}
\newtheorem{conclusion}[theorem]{Conclusion}
\newtheorem{condition}[theorem]{Condition}
\newtheorem{conjecture}[theorem]{Conjecture}
\newtheorem{corollary}[theorem]{Corollary}
\newtheorem{criterion}[theorem]{Criterion}
\newtheorem{definition}[theorem]{Definition}
\newtheorem{example}[theorem]{Example}
\newtheorem{exercise}[theorem]{Exercise}
\newtheorem{lemma}[theorem]{Lemma}
\newtheorem{notation}[theorem]{Notation}
\newtheorem{problem}[theorem]{Problem}
\newtheorem{proposition}[theorem]{Proposition}
\newtheorem{remark}[theorem]{Remark}
\newtheorem{solution}[theorem]{Solution}
\newtheorem{summary}[theorem]{Summary}
\newenvironment{proof}[1][Proof]{\noindent\textbf{#1.} }{\ \rule{0.5em}{0.5em}}
\begin{document}
\preprint{HEP/123-qed}
\title[Short title for running header]{REV\TeX4}
\author{First Author}
\affiliation{My Institution}
\author{Second Author}
\affiliation{My Institution}
\author{Third Author}
\affiliation{Other Institution}
\keywords{one two three}
\pacs{PACS number}

\begin{abstract}
Shell document for REV\TeX{} 4.

\end{abstract}
\volumeyear{year}
\volumenumber{number}
\issuenumber{number}
\eid{identifier}
\date[Date text]{date}
\received[Received text]{date}

\revised[Revised text]{date}

\accepted[Accepted text]{date}

\published[Published text]{date}

\startpage{101}
\endpage{102}
\maketitle
\tableofcontents


\section{Lattice}

A lattice is a partially ordered set in which every two elements have a
\emph{unique supremum, }$\vee,$ (also called a \emph{least upper bound} (LUB)
or \emph{join}) and a \emph{unique infimum, }$\wedge,$ (also called a
\emph{greatest lower bound} (GLB) or \emph{meet}).

A \emph{complemented lattice} is a \emph{bounded lattice} ($\equiv$\emph{with
least element 0 and greatest element 1}), in which every element $a$ has a
complement, i.e. an element $b$ satisfying
\begin{align}
a\vee b &  =1\\
a\wedge b &  =0.
\end{align}


\begin{remark}
Complements need not be unique.
\end{remark}

A \emph{relatively complemented lattice} is a lattice such that every
interval, $[c,d]$ viewed as a bounded lattice in its own right, is a
complemented lattice.

An \emph{orthocomplementation }on a complemented lattice is an
\emph{involution }which is \emph{order-reversing} and maps each element to a
complement. \emph{An orthocomplemented lattice satisfying a weak form of the
modular law is called an }\underline{\emph{orthomodular lattice.}}

In \emph{distributive lattices, complements are unique.} \emph{Every
complemented distributive lattice has a unique orthocomplementation and is in
fact a Boolean algebra.}

In general an element may have more than one complement. However, in a
(bounded) distributive lattice every element will have at most one complement.
A lattice in which every element has exactly one complement is called a
\emph{uniquely complemented} lattice

A lattice with the property that every interval (viewed as a sublattice) is
complemented is called a \emph{relatively complemented lattice}. In other
words, a relatively complemented lattice is characterized by the property that
for every element $a$ in an interval $[c,d]$ there is an element $b$ such
that
\begin{align}
a\vee b  &  =d\\
a\wedge b  &  =c.
\end{align}


Such an element $b$ is called a \emph{complement} of $a$ relative to the interval.

\emph{A distributive lattice is complemented if and only if it is bounded and
relatively complemented.}

\begin{remark}
\emph{The lattice of subspaces of a vector space provide an example of a
complemented lattice that is not, in general, distributive.}
\end{remark}

\subsection{Orthocomplementation}

An \emph{orthocomplementation on a bounded lattice} is a function that maps
each element $a$ to an "\emph{orthocomplement}" $a^{\bot}$ in such a way that
the following axioms are satisfied

\begin{enumerate}
\item Complement law%
\begin{align}
a^{\bot}\vee a  &  =1\\
a^{\bot}\wedge a  &  =0
\end{align}


\item Involution law%
\begin{equation}
\left(  a^{\bot}\right)  ^{\bot}=a
\end{equation}


\item Order-reversing%
\begin{equation}
a\leq b\Rightarrow b^{\bot}\leq a^{\bot}%
\end{equation}

\end{enumerate}

An \emph{orthocomplemented lattice or ortholattice is a bounded lattice which
is equipped with an orthocomplementation}. \emph{The lattice of subspaces of
an inner product space, and the orthogonal complement operation, provides an
example of an orthocomplemented lattice that is not, in general,
distributive.}

\emph{Boolean algebras are a special case of orthocomplemented lattices},
which in turn are a special case of \emph{complemented lattices} (with extra structure).

\begin{remark}
The closed subspaces of a separable Hilbert space represent quantum
propositions and behave as an orthocomplemented lattice. The
\emph{intersection} of two subspaces correspond to the \emph{meet} operation
and the \emph{union} of two subspaces to the \emph{join} operation i.e.
\begin{align}
V_{1}\cap V_{2}  &  \leftrightarrow V_{1}\wedge V_{2}\\
V_{1}\cup V_{2}  &  \leftrightarrow V_{1}\vee V_{2}%
\end{align}
and \emph{orthocomplementation} to \emph{orthogonality i.e.} the symbol
maintains in its meaning in the subspace and lattice settings.
\end{remark}

\begin{remark}
The \emph{intersection} of two subspaces is the largest subspace containd in
both of the subspaces, while the union of two subspaces is the smallest
subspace containing both the subspaces.
\end{remark}

\begin{remark}
The partial order is subspace inclusion i.e.
\begin{equation}
V_{n}\supset V_{n-1}\supset\cdots\supset V_{1}.
\end{equation}
Of course the index set does not have to be the natural numbers.
\end{remark}

\emph{Orthocomplemented lattices, like Boolean algebras, satisfy de Morgan's
laws:}%
\begin{align}
(a\vee b)^{\perp}  &  =a^{\bot}\wedge b^{\bot}\\
(a\wedge b)^{\perp}  &  =a^{\bot}\vee b^{\bot}%
\end{align}


\subsection{Orthomodular lattices}

A \emph{lattice }is called \emph{modular} if for all elements $a,b$ and $c$
\begin{equation}
a\leq c\Rightarrow a\vee(b\wedge c)=(a\vee b)\wedge c.
\end{equation}


This is \emph{weaker than distributivity}. A natural further weakening of this
condition for \emph{orthocomplemented lattices, (necessary for applications in
quantum logic)}, is to require it only in the special case
\begin{equation}
b=a^{\bot}.
\end{equation}
An \emph{orthomodular lattice} is therefore defined as an
\emph{orthocomplemented lattice} such that for any two elements
\begin{equation}
a\leq c\Rightarrow a\vee(a^{\bot}\wedge c)=(a\vee a^{\bot})\wedge c=c.
\end{equation}


If $\mathcal{A}$ is an arbitrary von Neumann algebra, then the set
$P(\mathcal{A})$ \emph{of all its projections is a complete orthomodular
lattice}. In these conditions, if $\mathcal{A}$ is a \emph{factor}, then a
\emph{dimension function }can be defined on the set $P(\mathcal{A})$.
Dependent on the set of values of this function, the factors are divided into
the types $I_{n},I_{\infty},II_{1},II_{\infty}$ and $III$ (the Murray--von
Neumann classification). It has been established that lattices of projections
of factors of the type $I_{n}$ and $II_{1}$ are \emph{continuous geometries.}

\begin{definition}
A continuous geometry is a lattice $\mathcal{L}$ with the following properties:
\end{definition}

\begin{enumerate}
\item $\mathcal{L}$ is modular.

\item $\mathcal{L}$ is complete.

\item The lattice operations $\wedge,\vee$ satisfy the continuity property%
\begin{equation}
\left(
%TCIMACRO{\tbigwedge \limits_{\alpha\in\mathfrak{A}}}%
%BeginExpansion
{\textstyle\bigwedge\limits_{\alpha\in\mathfrak{A}}}
%EndExpansion
a_{\alpha}\right)  \vee b=%
%TCIMACRO{\tbigwedge \limits_{\alpha\in\mathfrak{A}}}%
%BeginExpansion
{\textstyle\bigwedge\limits_{\alpha\in\mathfrak{A}}}
%EndExpansion
\left(  a_{\alpha}\vee b\right)  ,
\end{equation}
where $\mathfrak{A}$ is a directed set i.e. if $\U{3b1} <\U{3b2} $ then
$a_{\U{3b1} }<a\U{1d66} $,

\item The same condition with $\wedge$ and $\vee$ reversed i.e.
\begin{equation}
\left(
%TCIMACRO{\tbigvee \limits_{\alpha\in\mathfrak{A}}}%
%BeginExpansion
{\textstyle\bigvee\limits_{\alpha\in\mathfrak{A}}}
%EndExpansion
a_{\alpha}\right)  \wedge b=%
%TCIMACRO{\tbigvee \limits_{\alpha\in\mathfrak{A}}}%
%BeginExpansion
{\textstyle\bigvee\limits_{\alpha\in\mathfrak{A}}}
%EndExpansion
\left(  a_{\alpha}\wedge b\right)  ,
\end{equation}
.

\item Every element in $\mathcal{L}$ has a complement (not necessarily unique).

\item $\mathcal{L}$ is irreducible i.e. the only elements with unique
complements are $0$ and $1$.
\end{enumerate}

\begin{example}
Finite-dimensional complex projective space, or rather its set of linear
subspaces, is a continuous geometry,with dimensions taking values in the
discrete set $\{0,1/n,2/n,...,1\}$
\end{example}

\begin{example}
The projections of a finite type II$_{1}$ von Neumann algebra form a
continuous geometry with dimensions taking values in the unit interval $[0,1]$.
\end{example}

\begin{example}
Kaplansky (1955) showed that \emph{any} orthocomplemented complete modular
lattice is a continuous geometry.
\end{example}

\begin{example}
If $V$ is a vector space over a field (or division ring) $\mathbb{F}$, then
there is a natural map from the lattice $\mathcal{L}(V)$ of subspaces of $V$
to the lattice of subspaces of $V\otimes\mathbb{F}^{2}$ that multiplies
dimensions by $2$. So we can take a direct limit of%
\begin{equation}
PG(\mathbb{F})\subset PG(\mathbb{F}^{2})\subset PG(\mathbb{F}^{4})\subset
PG(\mathbb{F}^{8})\subset\cdots
\end{equation}
This has a dimension function taking values in all dyadic rationals between
$0$ and $1$. Its completion is a continuous geometry containing elements of
every dimension in $[0,1]$. This geometry was constructed by von Neumann and
is called the continuous geometry over $\mathbb{F}.$
\end{example}

\section{Projectors}

The lattice structure of the set of subspaces can be translated to algebraic
operations on the set of orthogonal projectors $P\left(  \mathcal{H}\right)
\subset\mathcal{B}\left(  \mathcal{H}\right)  ,$ the partial order
"$\supseteq"$ of subspace inclusion maps to "$\geq"$ of operator ordering and
orthocomplementation "$\bot"$ maps to orthogonality i.e.
\begin{equation}
P^{\bot}=I_{\mathcal{H}}-P.
\end{equation}
The meet and join operations follow from the use of the SVD\ and psuedo
inverse defined by
\begin{equation}
A^{+}=\sum_{1\leq j\leq\kappa}\left\vert e_{j}\right\rangle \left\vert
\alpha_{j}\right\vert ^{-1}\left\langle f_{j}\right\vert .
\end{equation}


\begin{theorem}
Let $A$ be an $m\times n$ matrix, and $B$ an $m\times k$ matrix both defined
on $\mathbb{C}$, then%
\begin{equation}
\operatorname{Ran}\left\{  AA^{\dagger}+BB^{\dagger}\right\}
=\operatorname{Ran}\left\{  A\right\}  \cup\operatorname{Ran}\left\{
B\right\}
\end{equation}

\end{theorem}

and%
\begin{equation}
\operatorname{Ker}\left\{  AA^{\dagger}+BB^{\dagger}\right\}
=\operatorname{Ker}\left\{  A\right\}  \cap\operatorname{Ker}\left\{
B\right\}  .
\end{equation}


\begin{corollary}
If $P$ and $Q$ are projections on $\mathbb{C}^{n}$
\begin{equation}
\operatorname{Ran}\left\{  P+Q\right\}  =\operatorname{Ran}\left\{  P\right\}
\cup\operatorname{Ran}\left\{  Q\right\}
\end{equation}
and
\begin{equation}
\operatorname{Ker}\left\{  P+Q\right\}  =\operatorname{Ker}\left\{  P\right\}
\cap\operatorname{Ker}\left\{  Q\right\}  .
\end{equation}

\end{corollary}

\begin{remark}
The sum of two projectors is not in general a projector i.e.%
\begin{equation}
\sum_{1\leq j\leq m}\left\vert e_{j}\right\rangle \left\langle e_{j}%
\right\vert +\sum_{1\leq j\leq n}\left\vert f_{j}\right\rangle \left\langle
f_{j}\right\vert \neq\sum_{1\leq j\leq\kappa}\left\vert h_{j}\right\rangle
\left\langle h_{j}\right\vert ,
\end{equation}
unless
\begin{equation}
\sum_{1\leq k\leq m}\sum_{1\leq j\leq n}\left\vert e_{j}\right\rangle
\left\langle e_{j}\right\vert \left\vert f_{k}\right\rangle \left\langle
f_{k}\right\vert =-\sum_{1\leq k\leq m}\sum_{1\leq j\leq n}\left\vert
f_{k}\right\rangle \left\langle f_{k}\right\vert \left\vert e_{j}\right\rangle
\left\langle e_{j}\right\vert
\end{equation}
i.e.%
\begin{equation}
\left\{  P,Q\right\}  =0.
\end{equation}

\end{remark}

\begin{theorem}
Let $M$ and $N$ be subspaces of $\mathbb{C}^{n}$. Let $P=P_{M}$ and $Q=P_{N}$.
Then all the following expressions are equal to the orthogonal projection
$P_{P\cup Q}$ on $M\cup N$ i.e.%
\begin{align}
P_{M\cup N}  &  =\left(  P+Q\right)  \left(  P+Q\right)  ^{+}\\
&  =\left(  P+Q\right)  ^{+}\left(  P+Q\right) \\
&  =Q+\left(  PQ^{\bot}\right)  ^{+}\left(  PQ^{\bot}\right) \\
&  =P+P^{\bot}\left(  P^{\bot}Q\right)  ^{+}\\
&  =Q+Q^{\bot}\left(  Q^{\bot}P\right)  ^{+}.
\end{align}

\end{theorem}

\begin{theorem}
Let $M$ and $N$ be subspaces of $\mathbb{C}^{n}$. Let $P=P_{M}$ and $Q=P_{N}$.
Then all the following expressions are equal to the orthogonal projection
$P_{P\cap Q}$ onto $M\cap N$ i.e.,
\begin{align}
P_{M\cap N}  &  =2Q\left(  P+Q\right)  P^{+}\\
&  =2P\left(  P+Q\right)  Q^{+}\\
&  =2\left(  P-P\left(  P+Q\right)  ^{+}P\right) \\
&  =P-\left(  P-QP\right)  ^{+}\left(  P-QP\right)  =P-\left(  Q^{\bot
}P\right)  ^{+}\left(  Q^{\bot}P\right) \\
&  =Q-\left(  Q-PQ\right)  ^{+}\left(  Q-PQ\right)  =Q-\left(  P^{\bot
}Q\right)  ^{+}\left(  P^{\bot}Q\right) \\
&  =P-P\left(  PQ^{\bot}\right)  ^{+}\\
&  =Q-Q\left(  QP^{\bot}\right)  ^{+}.
\end{align}


\begin{remark}
Von Neumann also showed that%
\begin{equation}
P_{M\cap N}=\lim\limits_{n\rightarrow\infty}\left\{  P\left(  QP\right)
^{n}\right\}  .
\end{equation}

\end{remark}
\end{theorem}

\section{Projectors in a Geometric Algebra}

The Geometric/Clifford algebra $\mathcal{F}$ can be viewed as the $w^{\ast}%
$-algebra $\mathcal{B}\left(  \mathcal{H}_{\mathcal{F}}\right)  ,$ where
$\mathcal{H}_{\mathcal{F}}$ is the exterior algebra, $%
%TCIMACRO{\tbigwedge }%
%BeginExpansion
{\textstyle\bigwedge}
%EndExpansion
\left(  \mathcal{H}^{1}\right)  $ of the Hilbert space $\mathcal{H}^{1}$ as
well as exterior algebra $%
%TCIMACRO{\tbigwedge }%
%BeginExpansion
{\textstyle\bigwedge}
%EndExpansion
\left(  \mathcal{F}_{1}\right)  ,$ that is equipped with a Clifford/Geometric product.

Consider the self adjoint conb $\left\{  \left.  \widetilde{x}_{j}\right\vert
1\leq j\leq2r\right\}  $ of the complex Geometric/Clifford algebra
$\mathcal{A}_{2r},$ where the a positive definite inner product in
$\mathcal{A}_{2r}$ is given by
\begin{equation}
\left\langle \left.  A\right\vert B\right\rangle _{\mathcal{F}}%
=\operatorname*{tr}\left\{  A^{\dagger}B\right\}  \equiv\left\langle
\left\langle A^{\dagger}B\right\rangle \right\rangle _{0},
\end{equation}
which also defines a norm
\begin{equation}
\left\Vert A\right\Vert _{\mathcal{F}}=\left(  \operatorname*{tr}\left\{
A^{\dagger}A\right\}  \right)  ^{\frac{1}{2}},
\end{equation}
thus making $\mathcal{F}$ a $c^{\ast}$, $w^{\ast}$ and \emph{Hilbert algebra},
when $r<\infty.$

\begin{remark}%
\begin{equation}
\left\langle A\right\rangle _{0}=\operatorname*{tr}\left\{  A\right\}  I_{2r}.
\end{equation}

\end{remark}

The psuedo scalar $\mathcal{J}_{2r}$ is given by%
\begin{equation}
\mathcal{J}_{2r}=%
%TCIMACRO{\tprod \limits_{1\leq j\leq2r}}%
%BeginExpansion
{\textstyle\prod\limits_{1\leq j\leq2r}}
%EndExpansion
\widetilde{x}_{j}=\widetilde{x}_{1}\cdots\widetilde{x}_{2r}%
\end{equation}
and has the properties

\begin{enumerate}
\item
\begin{align}
\mathcal{J}_{2r}\mathcal{J}_{2r}^{\dagger}  &  =I_{2r}\\
&  \Downarrow\\
\mathcal{J}_{2r}^{-1}  &  =\mathcal{J}_{2r}^{\dagger}%
\end{align}


\item
\begin{equation}
\left(  \mathcal{J}_{2r}\right)  ^{2}=-I_{2r}.
\end{equation}

\end{enumerate}

Hence
\begin{equation}
\mathcal{J}_{2r}^{-1}=\left(
%TCIMACRO{\tprod \limits_{1\leq j\leq2r}}%
%BeginExpansion
{\textstyle\prod\limits_{1\leq j\leq2r}}
%EndExpansion
\widetilde{x}_{j}\right)  ^{\dagger}=\widetilde{x}_{2r}\cdots\widetilde{x}%
_{1}.
\end{equation}


\subsection{Hodge Duality}

The \emph{dual} $A^{\%}$ of an $n$ -blade $A$ is defined by
\begin{equation}
A^{\%}=A\mathcal{J}_{2r}=\left(  -1\right)  ^{n\left(  r-n\right)
}\mathcal{J}_{2r}A.
\end{equation}
The action of the dual operator "$\%"$ on the basis $\left\{  \left.
\widetilde{x}_{j_{1}}\cdots\widetilde{x}_{j_{n}}\right\vert 1\leq j_{1}<\cdots
j_{n}\leq2r;0\leq n\leq2r\right\}  $ is
\begin{equation}
\left(  \widetilde{x}_{j_{1}}\cdots\widetilde{x}_{j_{n}}\right)
^{\%}=\widetilde{x}_{j_{1}}\cdots\widetilde{x}_{j_{n}}\widetilde{x}_{2r}%
\cdots\widetilde{x}_{1}\in\mathcal{A}_{2r-n},
\end{equation}
i.e.%
\begin{equation}
\%:\mathcal{F}_{n}\rightarrow\mathcal{F}_{2r-n}.
\end{equation}


\begin{example}
$r=2$
\begin{align}
I_{2r}\left(  \widetilde{x}_{4}\widetilde{x}_{3}\widetilde{x}_{2}%
\widetilde{x}_{1}\right)   &  =\widetilde{x}_{4}\widetilde{x}_{3}%
\widetilde{x}_{2}\widetilde{x}_{1}=\mathcal{J}_{4}^{-1}\\
\widetilde{x}_{1}\left(  \widetilde{x}_{4}\widetilde{x}_{3}\widetilde{x}%
_{2}\widetilde{x}_{1}\right)   &  =-\widetilde{x}_{4}\widetilde{x}%
_{3}\widetilde{x}_{2}\\
\widetilde{x}_{2}\left(  \widetilde{x}_{4}\widetilde{x}_{3}\widetilde{x}%
_{2}\widetilde{x}_{1}\right)   &  =\widetilde{x}_{4}\widetilde{x}%
_{3}\widetilde{x}_{1}\\
\widetilde{x}_{3}\left(  \widetilde{x}_{4}\widetilde{x}_{3}\widetilde{x}%
_{2}\widetilde{x}_{1}\right)   &  =-\widetilde{x}_{4}\widetilde{x}%
_{2}\widetilde{x}_{1}\\
\widetilde{x}_{4}\left(  \widetilde{x}_{4}\widetilde{x}_{3}\widetilde{x}%
_{2}\widetilde{x}_{1}\right)   &  =\widetilde{x}_{3}\widetilde{x}%
_{2}\widetilde{x}_{1}\\
\widetilde{x}_{1}\widetilde{x}_{2}\left(  \widetilde{x}_{4}\widetilde{x}%
_{3}\widetilde{x}_{2}\widetilde{x}_{1}\right)   &  =\widetilde{x}_{1}\left(
\widetilde{x}_{4}\widetilde{x}_{3}\widetilde{x}_{1}\right)  =\widetilde{x}%
_{4}\widetilde{x}_{3}%
\end{align}
The inner (geometric) and outer (wedge) products are related by duality.
Specifically, for $A_{n}=\left\langle A\right\rangle _{n}$ and $B_{m}%
=\left\langle B\right\rangle _{m}$%
\begin{equation}
A_{n}(B_{m}\mathcal{J}_{2r})=(A_{n}\wedge B_{m})\mathcal{J}_{2r}=\left(
-1\right)  ^{m\left(  2r-m\right)  }(A_{n}\mathcal{J}_{2r})B_{m}%
\end{equation}
or equivalently%
\begin{equation}
(A_{n}\wedge B_{m})^{\%}=A_{n}^{\%}B_{m}=\left(  -1\right)  ^{m\left(
2r-m\right)  }A_{n}B_{m}^{\%}.
\end{equation}
$\lceil$%
\begin{align}
\widetilde{x}_{1}\widetilde{x}_{2}\left(  \widetilde{x}_{4}\widetilde{x}%
_{3}\widetilde{x}_{2}\widetilde{x}_{1}\right)   &  =\widetilde{x}%
_{4}\widetilde{x}_{3}\\
\widetilde{x}_{1}\left(  \widetilde{x}_{4}\widetilde{x}_{3}\widetilde{x}%
_{2}\widetilde{x}_{1}\right)  \widetilde{x}_{2}  &  =-\widetilde{x}%
_{4}\widetilde{x}_{3}\
\end{align}
$\rfloor$
\end{example}

Right multiplication, $R_{\mathcal{J}_{2r}^{\dagger}},$ by $\mathcal{J}_{2r}$
is a unitary operator on $\mathcal{F}$ as
\begin{equation}
\left\langle \left.  A^{\%}\right\vert B^{\%}\right\rangle _{\mathcal{F}%
}=\left\langle \left.  A\mathcal{J}_{2r}\right\vert B\mathcal{J}%
_{2r}\right\rangle _{\mathcal{F}}=\operatorname*{tr}\left\{  \left(
A\mathcal{J}_{2r}\right)  ^{\dagger}B\mathcal{J}_{2r}\right\}
=\operatorname*{tr}\left\{  \mathcal{J}_{2r}^{\dagger}A^{\dagger}%
B\mathcal{J}_{2r}\right\}  =\operatorname*{tr}\left\{  A^{\dagger}B\right\}
=\left\langle \left.  A\right\vert B\right\rangle _{\mathcal{F}}%
\end{equation}
and the adjoint action $\widehat{\mathcal{J}_{2r}}$ is a unitary $c^{\ast}%
$-algebra \emph{involution} as
\begin{equation}
\left\langle \left.  \mathcal{J}_{2r}A\mathcal{J}_{2r}^{\dagger}\right\vert
\mathcal{J}_{2r}B\mathcal{J}_{2r}^{\dagger}\right\rangle _{\mathcal{F}%
}=\operatorname*{tr}\left\{  \left(  \mathcal{J}_{2r}A\mathcal{J}%
_{2r}^{\dagger}\right)  ^{\dagger}\mathcal{J}_{2r}B\mathcal{J}_{2r}^{\dagger
}\right\}  =\operatorname*{tr}\left\{  \mathcal{J}_{2r}A^{\dagger}%
\mathcal{J}_{2r}^{\dagger}\mathcal{J}_{2r}B\mathcal{J}_{2r}^{\dagger}\right\}
=\operatorname*{tr}\left\{  A^{\dagger}B\right\}  =\left\langle \left.
A\right\vert B\right\rangle _{\mathcal{F}}.
\end{equation}
$\lceil$%
\begin{equation}
\widetilde{x}_{1}\widetilde{x}_{2}\widetilde{x}_{3}\widetilde{x}%
_{4}\widetilde{x}_{1}\widetilde{x}_{2}\widetilde{x}_{4}\widetilde{x}%
_{3}\widetilde{x}_{2}\widetilde{x}_{1}=\widetilde{x}_{1}\widetilde{x}%
_{2}\widetilde{x}_{3}\widetilde{x}_{4}\widetilde{x}_{4}\widetilde{x}%
_{3}=\widetilde{x}_{1}\widetilde{x}_{2}%
\end{equation}
$\rfloor$

$\lceil$%
\begin{equation}
R_{\mathcal{J}_{2r}^{\dagger}}(A\wedge B)=R_{\mathcal{J}_{2r}^{\dagger}%
}\left(  A\right)  B=AB^{\%}=\left(  -1\right)  ^{m\left(  2r-m\right)
}AR_{\mathcal{J}_{2r}^{\dagger}}\left(  B\right)  .
\end{equation}
$\rfloor$

The \emph{wedge product} is also known as the \emph{outer product or
progressive product. One can define another product called the regressive
product, }$\vee,$\emph{ by }%
\begin{equation}
\left(  A_{n}\vee B_{m}\right)  ^{\%}=A_{n}^{\%}\wedge B_{m}^{\%}.
\end{equation}


\subsection{The Algebra of Subspaces}

To every $n$-dimensional subspace $V_{n}$ in $\mathcal{F}$ there corresponds
an $n-$blade $A_{n}=\left\langle A\right\rangle _{n}$ such that $V_{n}$ is the
solution set of the equation%
\begin{equation}
x\wedge A_{n}=0.
\end{equation}
Let $A_{n}$ be called the blade of $V_{n}$ while $V_{n}$ is called the support
of $A_{n}$. Since $A_{n}$ has a definite orientation and magnitude $\left\Vert
A_{n}\right\Vert _{\mathcal{F}}$, it associates an oriented weight or measure
with $V_{n}$. Thus, every blade in $\mathcal{F}$ determines a unique weighted
subspace of $V_{n}.$

Let $\mathfrak{B}=\left\{  \left.  B_{n}\right\vert 0\leq n\leq2r\right\}  $
be the set of all blades in $\mathcal{F}$ including the "scalars" as
$0$-blades. This set is closed under the inner and outer products. Therefore,
under these operations $\mathfrak{B}$ can be regarded as an algebra (\emph{not
a linear algebra}) of weighted subspaces in $\left\{  \left.  V_{n}\right\vert
0\leq n\leq2r\right\}  $. Note that $\mathfrak{B}$ is not closed under
addition, though some of its subsets are. With addition all of $\mathcal{F}$
can be generated from $\left\{  \left.  B_{n}\right\vert 0\leq n\leq
2r\right\}  $. For blades $A_{n}$ and $B_{m}$, we define a new product
$A_{n}\vee B_{m}$ by
\begin{equation}
A_{n}\vee B_{m}=A_{n}^{\%}B_{m}.
\end{equation}
This is Grassmann's regressive product. With this definition, Grassmann's
entire algebra of extension has been perfectly imbedded in geometric algebra.
All the properties of ther egressive product follow immediately from the
properties of the inner product and duality. Thus, it can be shown that for
\begin{equation}
\text{Grade}\left(  A\right)  +\text{Grade}\left(  B\right)  >n
\end{equation}
the regressive product is associative and is related to the outer product by
the \textquotedblleft De Morgan\textquotedblright\ rule%
\begin{equation}
(A\vee B)^{\%}=A^{\%}\wedge B^{\%}%
\end{equation}



\end{document}